\documentclass[
noshowanswers,        % 显示答案
answercolor=red,    % 答案颜色
paper=a4,           % 纸张大小
sectioncolor=cyan,  % 控制section颜色
subsectioncolor=cyan,
exercisecolor=orange, % 练习题颜色
examplecolor=cyan,    % 例题颜色
bianshicolor=purple  % 变式题颜色
%twocolumn 
]{biscuits}

% ============== 双栏间距控制 ==============
\setlength{\columnsep}{3em}       
\setlength{\columnseprule}{0.4pt}  

\usepackage{amssymb}      % 提供 \square 命令
\usepackage{tcolorbox}    % 提供彩色文本框
\tcbuselibrary{breakable} % 允许文本框跨页

\newenvironment{introduction}{%
	\begin{tcolorbox}[
		colback=orange!5,       % 背景色浅灰
		colframe=orange!30,     % 边框色深灰
		title=内容提要,        % 标题
		fonttitle=\bfseries,  % 标题加粗
		left=2mm, right=2mm,
		top=2mm, bottom=2mm,
		breakable             % 允许跨页
		]
		\renewcommand{\labelitemi}{$\square$} % 将 itemize 标签改为方框
		\begin{itemize}
			\setlength{\itemsep}{2pt} % 调整条目间距
			\setlength{\parskip}{0pt}
		}{%
		\end{itemize}
	\end{tcolorbox}
}

\usepackage{fontawesome5}   
\newtcolorbox{knowledge}[2][]{
	enhanced,
	breakable,
	arc=6pt,                         
	boxrule=0.8pt,                   
	colback=white,                   
	colframe=teal!70!black,          
	colbacktitle=teal!20,            
	coltitle=teal!80!black,          
	fonttitle=\sffamily\bfseries,    
	title={\faBook\ #2},             
	attach boxed title to top left={yshift=-3mm,xshift=5mm},
	boxed title style={
		boxrule=0pt,
		colback=teal!20,
		arc=4pt
	},
	drop fuzzy shadow=teal!30,       
	left=6mm, right=6mm, top=6mm, bottom=6mm,
	#1
}
\begin{document}
	
	% ========== 标题 ==========
	\title{ 一份适合高中数学题目排版的\LaTeX}
	\author{Biscuits}
	\date{2025.8}
	\maketitle
	
	\tableofcontents
	
	\newpage
	\chapter{第一章}
\section{导数的综合应用习题课}
\subsection{知识概要}
\begin{introduction}
	\item 掌握利用导数求函数单调区间、极值与最值的方法
	\item 理解导数与函数图像之间的关系，能绘制函数草图
	\item 熟练应用导数研究不等式恒成立与零点问题
\end{introduction}
\begin{knowledge}{导数应用核心步骤}
\parttitle{知识回顾}
	求解函数单调性、极值与最值问题需遵循标准化解题流程：第一步，优先确定函数定义域，所有导数分析、零点求解、极值判定均需在定义域范围内进行，杜绝定义域外无效解。第二步，对函数求导并彻底化简，结合四则求导、复合函数求导法则整理导函数，尽量化为因式乘积、分式最简形式，便于后续符号判断。第三步，求解临界点，令导函数为零求解驻点，同时找出定义域内导数不存在但函数连续的特殊点，统一作为可疑极值点。第四步，以临界点为分界点划分单调区间，列表判定各区间内导函数正负，确定函数增减性。第五步，根据导数左右符号变化判定极值，左增右减为极大值点，左减右增为极小值点，符号无变化则无极值。第六步，结合区间端点、极值点函数值，综合对比得到函数最值。

    
\end{knowledge}

	\subsection{题型精讲}
  \tixing{单调性}  
\begin{example}{1}[2025][全国][高三][模拟题][第1题]
	函数$f(x)=x^3-3x$的单调递增区间为\blankbox
	\begin{choices}
		\choice{$(-1,1)$}
		\choice{$(-\infty,-1),(1,+\infty)$}
		\choice{$(-\infty,-1)$}
		\choice{$(1,+\infty)$}
	\end{choices}
	\begin{proof}
		\begin{answer}
			B
		\end{answer}
		\begin{solutions}
			求导得$f'(x)=3x^2-3=3(x+1)(x-1)$，
			令$f'(x)>0$，解得$x<-1$或$x>1$，
			故单调递增区间为$(-\infty,-1),(1,+\infty)$，故选B。
		\end{solutions}
	\end{proof}
\end{example}

\begin{exercise}{1}[2025][全国][高三][模拟题][第2题]
	函数$f(x)=xe^{-x}$的极大值点为\blankbox
	\begin{choices}
		\choice{$x=-1$}
		\choice{$x=0$}
		\choice{$x=1$}
		\choice{$x=2$}
	\end{choices}
	\begin{proof}
		\begin{answer}
			C
		\end{answer}
		\begin{solutions}
			$f'(x)=e^{-x}-xe^{-x}=e^{-x}(1-x)$，
			当$x<1$时，$f'(x)>0$，$f(x)$单调递增；
			当$x>1$时，$f'(x)<0$，$f(x)$单调递减，
			故$x=1$为极大值点，选C。
		\end{solutions}
	\end{proof}
\end{exercise}

\begin{bianshi}{1}[2025][全国][高三][模拟题][第3题]
	函数$f(x)=x^2-2\ln x$的单调递减区间为\blankbox
	\begin{choices}
		\choice{$(-\infty,1)$}
		\choice{$(0,1)$}
		\choice{$(1,+\infty)$}
		\choice{$(0,+\infty)$}
	\end{choices}
	\begin{proof}
		\begin{answer}
			B
		\end{answer}
		\begin{solutions}
			函数定义域为$(0,+\infty)$，$f'(x)=2x-\dfrac{2}{x}=\dfrac{2(x^2-1)}{x}$，
			令$f'(x)<0$，结合$x>0$得$0<x<1$，
			故单调递减区间为$(0,1)$，选B。
		\end{solutions}
	\end{proof}
\end{bianshi}


\begin{example}{2}[2025][全国][高三][模拟题][第4题][多选]
	关于函数$f(x)=x^3-3x$，下列说法正确的是\blankbox
	\begin{choices}
		\choice{极大值为$2$}
		\choice{极小值为$-2$}
		\choice{单调递减区间为$(-1,1)$}
		\choice{在$\mathbb{R}$上单调递增}
	\end{choices}
	\begin{proof}
		\begin{answer}
			ABC
		\end{answer}
		\begin{solutions}
			$f'(x)=3(x+1)(x-1)$，
			$x\in(-\infty,-1)$时$f(x)$递增，$x\in(-1,1)$时$f(x)$递减，$x\in(1,+\infty)$时$f(x)$递增；
			$x=-1$时取极大值$f(-1)=2$，$x=1$时取极小值$f(1)=-2$，故ABC正确。
		\end{solutions}
	\end{proof}
\end{example}

\begin{example}{2}[2025][全国][高三][模拟题][第5题][多选]
	已知$f(x)=e^x-x$，则下列正确的是\blankbox
	\begin{choices}
		\choice{单调递减区间为$(-\infty,0)$}
		\choice{单调递增区间为$(0,+\infty)$}
		\choice{极小值为$1$}
		\choice{无极大值}
	\end{choices}
	\begin{proof}
		\begin{answer}
			ABCD
		\end{answer}
		\begin{solutions}
			$f'(x)=e^x-1$，令$f'(x)=0$得$x=0$；
			$x<0$时$f'(x)<0$，$f(x)$递减；$x>0$时$f'(x)>0$，$f(x)$递增；
			故$x=0$取极小值$f(0)=1$，无极大值，ABCD正确。
		\end{solutions}
	\end{proof}
\end{example}


\begin{example}{3}[2025][全国][高三][模拟题][第6题]
	函数$f(x)=x+\dfrac{1}{x}(x>0)$的极小值为\blankline
	\begin{proof}
		\begin{answer}
			$2$
		\end{answer}
		\begin{solutions}
			$f'(x)=1-\dfrac{1}{x^2}=\dfrac{x^2-1}{x^2}$，令$f'(x)=0$得$x=1$，
			$0<x<1$时$f(x)$递减，$x>1$时$f(x)$递增，
			故极小值$f(1)=1+1=2$。
		\end{solutions}
	\end{proof}
\end{example}

\begin{example}{3}[2025][全国][高三][模拟题][第7题]
	函数$f(x)=x^3-3x^2$的极大值为\blankline
	\begin{proof}
		\begin{answer}
			$0$
		\end{answer}
		\begin{solutions}
			$f'(x)=3x^2-6x=3x(x-2)$，令$f'(x)=0$得$x=0,x=2$，
			$x\in(-\infty,0)$递增，$(0,2)$递减，$(2,+\infty)$递增，
			极大值$f(0)=0$。
		\end{solutions}
	\end{proof}
\end{example}

\begin{example}{4}[2025][全国][高三][模拟题][第8题]
	已知函数$f(x)=x^3-3x^2+1$，求：
	\begin{enumerate}
		\item 函数$f(x)$的单调区间；
		\item 函数$f(x)$的极值。
	\end{enumerate}
	\begin{proof}
		\begin{answer}
			(1)单调递增区间$(-\infty,0),(2,+\infty)$，单调递减区间$(0,2)$；(2)极大值$1$，极小值$-3$
		\end{answer}
		\begin{solutions}
			$f'(x)=3x^2-6x=3x(x-2)$，令$f'(x)=0$，得$x=0,x=2$。
			当$x<0$或$x>2$时，$f'(x)>0$，$f(x)$单调递增；
			当$0<x<2$时，$f'(x)<0$，$f(x)$单调递减。
			故极大值$f(0)=1$，极小值$f(2)=-3$。
		\end{solutions}
	\end{proof}
\end{example}

\begin{bianshi}{4}[2025][全国][高三][模拟题][第22题]
	已知函数$f(x)=xe^{2x}-a(\ln x+x)$，其中$a\in\mathbb{R}$，$x>0$。
	\begin{enumerate}
		\item 讨论函数$f(x)$的单调性；
		\item 若函数$f(x)$存在两个不同极值点$x_1,x_2$，且$x_1<x_2$。
		\begin{enumerate}
			\item 证明：$2x_1+2x_2>-\ln a$；
			\item 求$a$的取值范围，并证明：$f(x_1)+f(x_2)>1+\ln2$。
		\end{enumerate}
	\end{enumerate}
	\begin{proof}
		\begin{answer}
			(1) 当$a\leq 2$时，$f(x)$在$(0,+\infty)$上单调递增；当$a>2$时，$f(x)$在$(0,x_1)$单调递增，在$(x_1,x_2)$单调递减，在$(x_2,+\infty)$单调递增；
			(2)(i) 不等式得证；(ii) $a\in(2,+\infty)$，不等式得证。
		\end{answer}
		\begin{solutions}
			\begin{enumerate}
				\item 函数定义域为$(0,+\infty)$，求导得$f'(x)=(2x+1)e^{2x}-a\left(\dfrac{1}{x}+1\right)=\dfrac{x(2x+1)e^{2x}-a(x+1)}{x}$。令$g(x)=x(2x+1)e^{2x}$，则$g'(x)=(4x^2+6x+1)e^{2x}$，$x>0$时$g'(x)>0$，故$g(x)$在$(0,+\infty)$单调递增。当$a\leq 2$时，$x(2x+1)e^{2x}>2(x+1)\geq a(x+1)$，$f'(x)>0$，$f(x)$在$(0,+\infty)$单调递增；当$a>2$时，存在$0<x_1<x_2$满足$g(x_1)=a(x_1+1),g(x_2)=a(x_2+1)$，故$f(x)$在$(0,x_1)$递增，$(x_1,x_2)$递减，$(x_2,+\infty)$递增。
				\item \begin{enumerate}
					\item 由$\begin{cases}x_1(2x_1+1)e^{2x_1}=a(x_1+1)\\x_2(2x_2+1)e^{2x_2}=a(x_2+1)\end{cases}$，两式取对数相加，结合$x_1x_2>\dfrac14$，整理放缩得$2x_1+2x_2>-\ln a$。
					\item 由(1)得$a\in(2,+\infty)$，代入化简得$f(x_1)+f(x_2)=2a-a\ln(x_1x_2)$，结合$x_1x_2>\dfrac14$与$a>2$，得$f(x_1)+f(x_2)>1+\ln2$。
				\end{enumerate}
			\end{enumerate}
		\end{solutions}
	\end{proof}
\end{bianshi}

\begin{example}{4}[2025][全国][高三][模拟题][第9题]
	已知函数$f(x)=x\ln x$，求$f(x)$的单调区间和极值。
	\begin{proof}
		\begin{answer}
			单调递减区间$\left(0,\dfrac{1}{e}\right)$，单调递增区间$\left(\dfrac{1}{e},+\infty\right)$；极小值$-\dfrac{1}{e}$，无极大值
		\end{answer}
		\begin{solutions}
			定义域$(0,+\infty)$，$f'(x)=\ln x+1$，令$f'(x)=0$得$x=\dfrac{1}{e}$。
			$0<x<\dfrac{1}{e}$时$f'(x)<0$，$f(x)$递减；$x>\dfrac{1}{e}$时$f'(x)>0$，$f(x)$递增。
			极小值$f\left(\dfrac{1}{e}\right)=-\dfrac{1}{e}$，无极大值。
		\end{solutions}
	\end{proof}
\end{example}
	
	\subsection{易错警示}
	\begin{knowledge}{导数题常见错误}
		1. 忘记求定义域；2. 求导公式错误；3. 极值点概念混淆（导数为零的点不一定是极值点）；4. 求最值时忽略比较区间端点值；5. 含参讨论忽略分类标准。
	\end{knowledge}
	
\end{document}